% Part: sets-functions-relations
% Chapter: functions
% Section: functions-relations

\documentclass[../../../include/open-logic-section]{subfiles}


\begin{document}

\olfileid{sfr}{fun}{rel}

\olsection{Functies als relaties}

\begin{explain}
Een functie die !!{element}s van~$A$ naar !!{element}s van~$B$ stuurt,
legt meteen een relatie tussen $A$ en~$B$ vast. Die relatie geldt voor
$x$ en $y$ precies wanneer $f(x) = y$. Soms willen we de wiskunde met
zo weinig mogelijk bouwstenen opbouwen. Dan kunnen we de functie~$f$
\emph{zien als} deze relatie, dus als een verzameling paren. De vraag
is dan: welke relaties leggen op deze manier een functie vast?
\end{explain}

\begin{defn}[Grafiek van een functie] Neem een functie $f\colon A \to
B$. De \emph{grafiek} van~$f$ is de relatie $R_f \subseteq A \times
B$ die als volgt is vastgelegd:
\[
R_f = \Setabs{\tuple{x,y}}{f(x) = y}.
\]
\end{defn}

\begin{explain}
Door extensionaliteit ligt de grafiek van een functie helemaal vast.
Dezelfde regel geeft ook de stilzwijgende regel van extensionaliteit
voor functies: hebben $f$ en~$g$ hetzelfde domein en codomein, dan
zijn ze gelijk als ze bij elke invoer dezelfde waarde geven.

Net zo geldt: is een relatie ‘functioneel’, dan is zij de grafiek van
een functie.
\end{explain}

\begin{prop}\ollabel{prop:graph-function}
Neem $R \subseteq A \times B$ waarvoor de volgende twee dingen gelden:
\begin{enumerate}
\item Als $Rxy$ en $Rxz$, dan $y = z$; en
\item bij elke $x \in A$ is er een $y \in B$ waarvoor $\tuple{x,
y} \in R$.
\end{enumerate}
Dan is $R$ de grafiek van de functie $f\colon A \to B$ die is
vastgelegd door $f(x) = y$ precies wanneer $Rxy$.
\end{prop}

\begin{proof}
Neem aan dat er een $y$ is waarvoor $Rxy$ geldt. Als er ook een $z
\neq y$ was waarvoor $Rxz$ geldt, zou $R$ niet aan de voorwaarde
voldoen. Als er een $y$ is waarvoor $Rxy$ geldt, is deze $y$ dus de
enige. Daarom is $f$ goed vastgelegd. Het is duidelijk dat $R_f = R$.
\end{proof}

\begin{explain}
Elke functie $f\colon A \to B$ heeft een grafiek: een relatie die uit
paren in $A \times B$ bestaat en is vastgelegd door $f(x) = y$.
Omgekeerd is elke
relatie~$R \subseteq A \times B$ met de eigenschappen uit
\olref{prop:graph-function} de grafiek van een functie~$f \colon A \to
B$. Door dit nauwe verband kunnen we een functie gewoon als haar
grafiek zien. Met andere woorden: we kunnen functies zien als bepaalde
relaties, dus als bepaalde verzamelingen tupels.
\oliflabeldef{sfr:rel:ref:sec}{Let wel: met deze ‘identificatie’
bedoelen we hetzelfde als in \olref[sfr][rel][ref]{sec}. We doen geen
uitspraak over wat functies in diepste zin zijn; we merken alleen op
dat het handig is om functies als bepaalde verzamelingen te
\emph{zien}. Dat heeft dit voordeel: w}{W}e kunnen nu op functies
dezelfde soort bewerkingen uitvoeren als eerder op relaties (zie
\olref[sfr][rel][ops]{sec}). In het bijzonder:
\end{explain}

\begin{defn}\ollabel{defn:funimage}
Neem een functie $f \colon A \to B$ met $C\subseteq A$.

De \emph{beperking} van~$f$ tot~$C$ is de
functie~$\funrestrictionto{f}{C}\colon C \to B$ waarvoor
$(\funrestrictionto{f}{C})(x) = f(x)$ voor alle $x \in C$. Anders
gezegd: $\funrestrictionto{f}{C} = \Setabs{\tuple{x, y} \in R_f}{x
\in C}$.

Als we~$f$ op~$C$ \emph{toepassen}, krijgen we $\funimage{f}{C} =
\Setabs{f(x)}{x \in C}$. We noemen dit ook het \emph{beeld} van~$C$
onder~$f$.
\end{defn}

\begin{explain}
Uit deze definities volgt dat $\ran{f} =
\funimage{f}{\dom{f}}$ voor elke functie~$f$.
\oliflabeldef{sfr:rel:ops:sec}{Deze begrippen werken precies zoals je
zou verwachten uit de definities in
\olref[sfr][rel][ops]{sec} en uit het feit dat we functies als relaties
zien. Twee andere bewerkingen hebben wat meer uitleg nodig: inversen nemen en
relaties samenstellen. Die geven we in
\olref[inv]{sec} en \olref[cmp]{sec}.}{}
\end{explain}
\end{document}
